\chapter{Scheduling Theory and Task Allocation}\label{chap:schedulingtheory}

With the motivation of modelling real-world mathematical development in mind, we are led to a general problem of allocating heterogeneous agents (i.e. mathematicians) across a large pool of interdependent preemptimble tasks (i.e. theorems). These agents operate in a decentralized manner, with tasks having a natural implicit dependency structure between them with agents able to generate subtasks (e.g., conjectures, lemmas) that may assist in the completion of other tasks. The goal is to maximize collective progress towards some target set of tasks (e.g., the proof of a major conjecture) in a minimal amount of time.

This problem of interdependent task allocation is closely related to classical scheduling theory, which studies how to allocate jobs to machines under various constraints and objectives. In the standard taxonomy of machine scheduling, one studies how jobs should be allocated to machines under objectives such as makespan or completion time, with increasingly general machine models ranging from identical to uniform to unrelated machines. 

More fundamentally, one may consider the problem of scheduling along three orthogonal axes: deterministic vs. stochastic environment, offline vs. online environment, and deterministic vs. stochastic policy. We review the relevant scheduling theory along these axes, noting that deterministic environments are a special case of the more general stochastic environment, and that offline environments are a special case of the more general online environment, and deterministic policies are a special case of the more general stochastic policies.

\subsection{Offline Deterministic Scheduling}
Classical scheduling theory provides a mathematical foundation for understanding how to optimally allocate jobs to machines under various constraints. Its importance lies in offering formal models and algorithms for resource allocation problems that arise ubiquitously in manufacturing, computing, project management, and logistics. By developing a taxonomy of machine environments, job characteristics, and objectives, scheduling theory enables the robust analysis and study of such resource allocation tasks.

More rigorously, given $n$ jobs $\{J_j\}_{j=1}^n$ and $m$ machines $\{M_i\}_{i=1}^m$, we may characterize offline deterministic scheduling problems by Graham, Lawler, Lenstra, and Rinnooy-Kan's three-field notation $\alpha \mid  \beta \mid  \gamma$, consisting of a machine environment $\alpha$, a job environment $\beta$, and an objective $\gamma$.

\begin{definition}[Machine Environment]
    A machine environment $\alpha=\alpha_1\alpha_2$, given by $\alpha_1 \in \{\epsilon,P,Q,R,O,F,J\}, \alpha_2 \in \{\epsilon\}\cup \mathbb{N}$ is defined as follows:
    \begin{itemize}
        \item $\alpha_1 \in \{\epsilon,P,Q,R\}$: each $J_j$ consists of a single operation processable on any $M_i$ with processing time $p_{ij}$ as follows:
        \begin{itemize}
            \item $\alpha_1 = \epsilon$: single machine $p_{1j}=p_j$
            \item $\alpha_1 = P$: identical parallel machines $p_{ij}=p_j$
            \item $\alpha_1 = Q$: uniform parallel machines $p_{ij} = p_j/q_i$ where $q_i$ is the speed of machine $M_i$
            \item $\alpha_1 = R$: unrelated parallel machines with arbitrary $p_{ij}$
        \end{itemize}
        \item $\alpha_1 = O$: open shop with operation set $\{O_{1j},\cdots,O_{mj}\}$ for each job $J_j$. $O_{ij}$ must be processed on machine $M_i$ in $p_{ij}$ time, but the order is arbitrary., each operation processable on any machine with processing time $p_{ij}$.
        \item $\alpha_1 = J$: job shop with operations $(O_{1j},\cdots,O_{m_j j})$ for each job $J_j$. $O_{ij}$ must be processed on machine $\mu_{ij}$ in $p_{ij}$ time, with $\mu_{i-1,j} \neq \mu_{ij}$ for $i = 2,\cdots,m_j$. 
        \item $\alpha_1 = F$: flow shop with operations $(O_{1j},\cdots,O_{mj})$ for each job $J_j$. $O_{ij}$ must be processed on machine $M_i$ in $p_{ij}$ time, and the order is fixed such that $O_{1j}$ must be processed before $O_{2j}$, which must be processed before $O_{3j}$, etc.
    \end{itemize}

    Additionally, if $\alpha_2 \in \mathbb{N}$, there are $m=\alpha_2$ machines; otherwise, $m$ is assumed to be variable. Note, $\alpha_1=\epsilon \iff \alpha_2=1$.
\end{definition}

\begin{definition}[Job Environment]
    A job environment $\beta \subseteq \{\beta_1,\cdots,\beta_6\}$ is given as a collection of job characteristics, which are given as follows:
    \begin{itemize}
        \item $\beta_1 \in \{\mathrm{pmtn}, \epsilon\}$: If $\beta_1=\mathrm{pmtn}$, preemption is allowed (an operation may be interrupted and resumed later). Otherwise, no preemption.

        \item $\beta_2 \in \{\mathrm{res}, \mathrm{res1}, \epsilon\}$: If $\beta_2=\mathrm{res}$, then $s$ limited resources $R_h$ are present; job $J_j$ requires $r_{hj}$ units of $R_h$ throughout its processing, and resource capacities cannot be exceeded. If $\beta_2=\mathrm{res1}$, a single limited resource is present and otherwise for $\beta_2=\epsilon$, there are no resource constraints.

        \item $\beta_3 \in \{\mathrm{prec}, \mathrm{tree}, \epsilon\}$: If $\beta_3=\mathrm{prec}$, then a precedence relation $\prec$ between jobs is defined via a DAG $G$ with vertex set $[n]$. If there exists a directed path from $j$ to $k$, then $J_j\prec J_k$ and $J_j$ must complete before $J_k$ starts.
        
        If $\beta_3=\mathrm{tree}$, we moreover require that $G$ is a rooted tree. Otherwise, if $\beta_3=\epsilon$, there are no precedence constraints.

        \item $\beta_4 \in \{r_j,\epsilon\}$: If $\beta_4=r_j$, job-specific release dates are specified, and otherwise, all jobs are available at time $0$.

        \item $\beta_5 \in \{m_j \le \bar{m},\epsilon\}$ (only relevant for $\alpha_1=J$): If $\beta_5=m_j \le \bar{m}$, then a constant upper bound on the number of operations per job is specified. Otherwise, no bound on operations is specified.
        
        \item $\beta_6 \in \{p_{ij}=1,\; p^- \le p_{ij} \le p^+,\; \epsilon\}$: If $\beta_6=p_{ij}=1$, all operation processing times are unit. Otherwise, if $\beta_6=p^- \le p_{ij} \le p^+$, uniform lower/upper bounds on processing times are specified. If $\beta_6=\epsilon$, we have no such bounds.
    \end{itemize}
\end{definition}

\begin{definition}[Objective]
    An objective $\gamma \in \{f_{\max}, f_{\sum}\}$ indicates the optimality criterion one is optimizing for. Namely, given a schedule, one may compute for each job $J_j$:
    \begin{itemize}
        \item total completion time $C_j$
        \item latency $L_j = C_j - d_j$
        \item tardiness $T_j = \max\{0, C_j - d_j\}$
        \item unit penalty $U_j = \begin{cases}
            1 & \text{if } C_j > d_j \\ 0 & \text{otherwise}
        \end{cases}$
    \end{itemize}
    Where the latter three are well defined if due dates $d_j$ are specified for each job. Then, if $\gamma = f_{\max}$, the objective is to minimize
    $f_{\max} \in \{C_{\max}, L_{\max}\}$, where $C_{\max} = \max_{j} C_j$ is the maximum completion time and $L_{\max} = \max_{j} L_j$ is the maximum latency.
    
    Alternatively, if $\gamma = f_{\sum}$, the objective is to minimize
    $$f_{\sum} \in \left\{\sum_j C_j,\; \sum_j T_j,\; \sum_j U_j,\; \sum_j w_j C_j,\; \sum_j w_j T_j,\; \sum_j w_j U_j\right\}$$
    for some weights $w_j$.
\end{definition}

With this, we may naturally define an instance of the (offline) deterministic scheduling problem as:

\begin{definition}[Offline Deterministic scheduling instance]
    An offline deterministic scheduling instance $I$ is given as a tuple:
    $$
        I = (\mathcal J, \mathcal M, \alpha \mid  \beta \mid  \gamma)
    $$
    Augmented with any additional data required by $\alpha, \beta,\gamma$, such as processing times $p_{ij}$, precedence DAG $G$, job operations $O_j$, resource requirements $r_{hj}$, objective weights $w_j$, etc.

    In practice, we say this is an (offline deterministic) instance of type $\alpha\mid \beta\mid \gamma$ with jobs $\mathcal{J} = \{J_1,\dots,J_n\}$ and machines $\mathcal{M} = \{M_1,\dots,M_m\}$, where the machine environment $\alpha$, job environment $\beta$, and objective $\gamma$ are given as above.
\end{definition}


\subsection{Offline Stochastic Scheduling}

Real-world environments are often subject to various sources of uncertainty, such as variability in processing times, machine breakdowns, etc. To model such phenomena, one may extend the classical scheduling framework to incorporate stochastic primitives, leading to the study of optimal scheduling under uncertainty.

Towards this, we consider scheduling models in which the distributions of processing times, release dates, due dates, or other relevant parameters are known at time zero, however, we note that the actual outcome of a random processing time is known only at the time of its completion. Similarly, the realization of a release date or due date is known only at the point in time when it occurs. 


\begin{definition}[Offline Stochastic scheduling instance]

An (offline) stochastic scheduling instance $I$ is given as a tuple:
$$
I= (\mathcal{J}, \mathcal{M}, \alpha \mid  \beta \mid  \gamma, (\Omega, \mathcal{F}, \mathbb{P}))
$$

Namely, we consider such an instance with job set $\mathcal{J} = \{J_1,\dots,J_n\}$ and machines $\mathcal{M} = \{M_1,\dots,M_m\}$ of type $\alpha \mid  \beta \mid  \gamma$ defined with respect to a probability space $(\Omega, \mathcal{F}, \mathbb{P})$. The machine environment $\alpha$, job environment $\beta$, and objective $\gamma$ are defined as in the deterministic case, but with primitives now allowed to be given as a random variable.

More rigorously, a stochastic scheduling instance of type $\alpha\mid \beta\mid \gamma$ is defined analogously to the deterministic case, but with the following modifications:
\begin{itemize}
    \item processing times $p_{ij}$ may be given as random variables $P_{ij}:\Omega\to(0,\infty)$
    \item if $\alpha_1=Q$, machine speeds $q_i$ may be given as random variables $Q_i:\Omega\to(0,\infty)$, and processing times are given by $P_{ij} = P_j / Q_i$.
    \item if $\alpha_1\in\{O,J,F\}$, for each job $J_j$, the processing time of each operation $O_{ij} \in \{O_{1j},\dots,O_{m_j j}\}$ may be given as a random variable $P_{ij}:\Omega\to (0,\infty)$. 
    \item If $\beta_2 = \mathrm{res}$, the resource requirements $r_{hj}$ for a job $J_j$ may be given as random variables $R_{hj}:\Omega\to[0,\infty)$.
    \item If $\beta_4 = r_j$, the release date of job $J_j$ may be given as random variables $R_j:\Omega\to[0,\infty)$.
    \item If $\gamma$ relies on due dates, the due date of job $J_j$ may be given as random variables $D_j:\Omega\to[0,\infty)$.
\end{itemize}

Additionally, we extend the objective \(\gamma\) to stochastic analogues of the aforementioned objectives; in practice, this most commonly occurs through expectation: $\Gamma \in \{\mathbb{E}(f_{\max}), \mathbb{E}(f_{\sum})\}$. For example, one may consider the expected makespan objective $\Gamma = \mathbb{E}(C_{\max})$.
\end{definition}
\begin{remark}[Latent realization versus observation]
All primitive random variables in an offline stochastic instance are defined on the same probability space. Hence, for each \(\omega\in\Omega\), the full collection of primitive parameters already has a realized value.
\end{remark}

\begin{remark}[Observation convention]
The bare instance definition specifies the latent random primitives, but not yet the scheduler's information process. In the standard stochastic scheduling convention, the distributions of the primitive random variables are known in advance, while their realized values are only observed when they become operationally relevant. Typical conventions are:
\begin{itemize}
    \item a realized processing requirement is observed upon completion of the corresponding job or operation;
    \item a realized release date is observed when the job is released;
    \item a realized due date is either known at time \(0\) or else observed when it becomes relevant, depending on the model;
    \item any other random primitive is observed according to the information convention specified later when defining states and policies.
\end{itemize}
\end{remark}


In this setting, the scheduler must make decisions based on the information available at each point in time, without knowledge of future random outcomes. Indeed, our objective in stochastic scheduling is not simply to design algorithms that optimize our desired objective $\gamma$, but rather to do so in a nonanticipative manner. This will be formalized in more rigor and generality in the following section.

\subsection{Stochastic Online Scheduling}

In online scheduling, the defining feature is not merely uncertainty, but sequential revelation of the instance. The scheduler does not know the full instance at time \(0\); instead, jobs, operations, precedence constraints, or other problem data are revealed over time, and decisions must be made using only the information currently available. In the fully stochastic regime, the unrevealed portion of the instance may itself be random. Here, we consider the more general stochastic online scheduling setting, which subsumes the deterministic online scheduling setting as a special case.


\begin{definition}[Stochastic online scheduling instance]
A stochastic online scheduling instance \(I\) is given by a tuple
\[
I=(\mathcal J^\ast,\mathcal M,\alpha\mid\beta\mid\gamma,(\Omega,\mathcal F,\mathbb P),\mathcal V),
\]

Namely, we consider an instance with a (perhaps countable) latent job set $\mathcal{J}^*$ and machine set $\mathcal{M} = \{M_1,\dots,M_m\}$ of type $\alpha\mid \beta\mid \gamma$ defined with respect to a probability space $(\Omega, \mathcal{F}, \mathbb{P})$. Here, we additionally include a family of revelation-time random variables $\mathcal{V}=(V_j)_{J_j\in\mathcal J^\ast}$, where $V_j:\Omega\to[0,\infty)$ denotes the time at which job $J_j$ becomes visible to the scheduler.



Note that if $\beta_3 = \mathrm{prec}$ or $\beta_3 = \mathrm{tree}$, the precedence DAG $G^* = (\mathcal{J}^*, E^*)$ is also latent and revealed over time as jobs are revealed. In particular, if we have the set of currently revealed jobs $\mathcal{J}_t(\omega) = \{J_j \in \mathcal{J}^* : V_j(\omega) \le t\}$, then the scheduler only observed precedence constraints between jobs in $\mathcal{J}_t(\omega)$ induced via the sub-DAG $G_t(\omega) = (\mathcal{J}_t(\omega), E_t(\omega))$ for $E_t(\omega) = \{(J_j, J_k) \in E^* : J_j, J_k \in \mathcal{J}_t(\omega)\}$.


Moreover, in this precedence-constrained setting, we note that an instance is valid only if the revelation times are consistent with the precedence constraints. Namely, if $J_j \prec J_k$, then we must have $V_j \le V_k$ almost surely.

As in the offline stochastic case, the numerical primitives attached to each latent job may be random variables on \((\Omega,\mathcal F,\mathbb P)\). In particular, consider the random variables:
\begin{itemize}
    \item processing requirements \(P_{ij}:\Omega\to(0,\infty)\) or \(P_j:\Omega\to(0,\infty)\);
    \item if \(\alpha_1=Q\), machine speeds \(Q_i:\Omega\to(0,\infty)\);
    \item if \(\alpha_1\in\{O,J,F\}\), operation processing requirements \(P_{ij}:\Omega\to(0,\infty)\);
    \item if \(\beta_2=\mathrm{res}\), resource requirements \(R_{hj}:\Omega\to[0,\infty)\);
    \item if \(\beta_4=r_j\), release dates \(R_j:\Omega\to[0,\infty)\). We distinguish between release date and revelation times for the sake of comprehensiveness, only requiring $R_j \ge V_j$. However, in practice, one often assumes $R_j = V_j$;
    \item if \(\gamma\) depends on due dates, due dates \(D_j:\Omega\to[0,\infty)\).
\end{itemize}

The defining feature of the online model is that the full latent job set \(\mathcal J^\ast\) and its associated primitive data need not be visible to the scheduler at time \(0\). Instead, a job \(J_j\) becomes available for consideration only when it is revealed at time \(V_j\).
\end{definition}

\begin{remark}[Revealed job set]
As noted above, For each \(t\ge 0\) and \(\omega\in\Omega\), we define the set of jobs visible to the scheduler by time \(t\) as
\[
\mathcal J_t(\omega):=\{J_j\in\mathcal J^\ast : V_j(\omega)\le t\}.
\]
Thus \((\mathcal J_t)_{t\ge 0}\) is the revealed-job process induced by the latent instance and the revelation times \(V_j\).
\end{remark}

\begin{remark}[Observation convention]
Unless otherwise specified, when a job \(J_j\) is revealed at time \(V_j\), the scheduler learns the existence of \(J_j\) together with any job-level metadata designated as visible upon revelation. In the simplest convention, this includes data such as due dates, release dates, resource requirements, or the distributions of latent processing requirements.

As in the offline stochastic case, the realized processing requirement itself need not be observed upon revelation. Instead, it may remain latent and become observable only when the corresponding job or operation completes.
\end{remark}

With this, we may simply define the deterministic online scheduling setting as a special case of the stochastic online scheduling setting:
\begin{definition}[Deterministic online scheduling instance]
A deterministic online scheduling instance is a stochastic online scheduling instance in which all primitives are deterministic, i.e. all random variables are degenerate.
\end{definition}

\subsection{States, Actions, and Policies}

With the scheduling instance now defined across the deterministic/stochastic and
offline/online axes, we turn to formalizing the scheduler's decision problem. The central
question is: given a stochastic online scheduling instance, what does it mean to make
decisions optimally, and what mathematical objects describe the scheduler's strategy?

In the offline deterministic setting, this question is trivial to formulate (though not to
solve): the scheduler knows the full instance at time $0$, and the problem reduces to
combinatorial optimization over the set of feasible schedules. No sequential
decision-making under uncertainty is involved.

Once stochasticity or online revelation is introduced, however, the problem acquires a
fundamentally sequential character. The scheduler must choose actions at each time step
based only on the information available so far, without knowledge of unrevealed jobs or
unrealized processing times. The natural mathematical framework for such problems is
the theory of partially observable Markov decision processes (POMDPs), in which a
latent Markov process evolves according to a transition kernel, and the controller
observes only a partial projection of the true state at each step.

We now recall the general POMDP framework, then show that the stochastic online
scheduling problem admits a natural POMDP formulation. The offline stochastic, online
deterministic, and offline deterministic settings will each be recovered as special cases.
We work throughout in discrete time $t \in \mathbb{N}_0$.

\subsubsection{Partially Observable Markov Decision Processes}

A partially observable Markov decision process models sequential decision-making in an
environment whose true state is hidden from the controller. At each time step, the
controller receives a partial observation of the state, selects an action, and the
environment transitions to a new state. The controller's goal is to select actions so as
to optimize a cumulative performance criterion, using only the history of observations
and actions available to it.

\begin{definition}[POMDP]
A discrete-time partially observable Markov decision process is a tuple
\[
\mathcal{P} = (\mathsf{S},\; \mathsf{A},\; \mathsf{O},\; T,\; T^0,\; Z,\; c),
\]

for a state space $\mathsf{S}$, a measurable space whose elements encode the full (hidden) state of the environment; an action space $\mathsf{A}$, a measurable space of actions available to the controller; an observation space $\mathsf{O}$, a measurable space of observations received by the controller; a transition kernel $T(\cdot \mid s, a) \in \mathcal{P}(\mathsf{S})$ specifying the distribution of the next state given a current state and action; an initial state distribution $T^0 \in \mathcal{P}(\mathsf{S})$; an observation kernel $Z(\cdot \mid s) \in \mathcal{P}(\mathsf{O})$ specifying the distribution of the observation received by the controller when the environment is in state $s$; and a per-step cost function $c : \mathsf{S} \times \mathsf{A} \to \mathbb{R}$.

\end{definition}

The process evolves as follows. At time $t = 0$, nature draws a latent state
$s_0 \sim T^0$ and the controller receives an observation $o_0 \sim Z(\cdot \mid s_0)$.
At each subsequent time $t \ge 0$, the controller selects an action $a_t$, the environment
transitions to a new state $s_{t+1} \sim T(\cdot \mid s_t, a_t)$, and the controller receives
a new observation $o_{t+1} \sim Z(\cdot \mid s_{t+1})$.

The controller's observation history at time $t$ is the tuple
\[
h_t := (o_0,\, a_0,\, o_1,\, a_1,\, \ldots,\, a_{t-1},\, o_t),
\]
which constitutes the controller's complete information record. A randomized nonanticipative policy is a family of stochastic kernels
$\pi = (\pi_t)_{t \ge 0}$ with $\pi_t : \mathsf{H}_t \to \mathcal{P}(\mathsf{A})$, where $\mathsf{H}_t$ is the space of observation histories of length $t$, such that the action $a_t \sim \pi_t(\cdot \mid h_t)$ depends only on the observation history up to time $t$. 

Equivalently, the action process $(a_t)_{t \ge 0}$ is adapted to the observation filtration $\mathcal{F}_t^{\mathrm{obs}} := \sigma(o_0, a_0, \ldots, a_{t-1}, o_t)$.

\begin{remark}[Belief state]
By standard POMDP theory, the posterior belief
$b_t := \mathrm{Law}(s_t \mid h_t) \in \mathcal{P}(\mathsf{S})$ is a sufficient statistic of
the observation history for optimal control: there exists an optimal policy whose action
at time $t$ depends on $h_t$ only through $b_t$. One may equivalently reformulate the
POMDP as a fully observed MDP on the belief space $\mathcal{P}(\mathsf{S})$, the belief MDP, in which the state is $b_t$ and the transition is given by Bayesian updating.
\end{remark}

When the state space, observation space, and action space carry additional structure, it is common to restrict the action space in a state-dependent or observation-dependent manner. If the controller's available actions depend only on the current observation (as will be the case in our scheduling formulation), one specifies an admissible
action correspondence $\mathcal{U} : \mathsf{O} \rightrightarrows \mathsf{A}$ and requires that $\pi_t(\mathcal{U}(o_t) \mid h_t) = 1$ for all $t$ and $h_t$.

We will formulate the stochastic online scheduling problem as a stochastic shortest-path POMDP: a POMDP equipped with a terminal set
$\mathsf{S}^{\mathrm{term}} \subseteq \mathsf{S}$ and per-step cost $c$, where the objective is to minimize the expected cumulative cost until termination,
\[
\inf_{\pi} \; \mathbb{E}^{\pi}\!\left[\sum_{t=0}^{\tau_{\mathrm{term}}-1}
c(s_t, a_t)\right],
\]
with $\tau_{\mathrm{term}} := \inf\{t \ge 0 : s_t \in \mathsf{S}^{\mathrm{term}}\}$.

For expected makespan, we will take $c \equiv 1$ until all target jobs complete, so that the cumulative cost equals $\tau_{\mathrm{term}} = C_{\max}$ and the objective reduces to $\inf_\pi \mathbb{E}^\pi[C_{\max}]$.

\subsubsection{Scheduling as a POMDP}

We now show that the stochastic online scheduling problem defined in the preceding sections admits a natural formulation as a stochastic shortest-path POMDP. Namely, we fix a stochastic online scheduling instance $I$ and assume there exists a finite target job set $\mathcal{J}^{\mathrm{target}} \subseteq \mathcal{J}^*$ such that the scheduling
objective is evaluated with respect to the completion of the jobs in $\mathcal{J}^{\mathrm{target}}$. Notably, if $\mathcal{J}^*$ is finite, one can simply take $\mathcal{J}^{\mathrm{target}} = \mathcal{J}^*$.

With this in mind, we now define the scheduling POMDP induced by a stochastic online scheduling instance.

\begin{definition}[Scheduling POMDP]\label{def:schedpomdp}
Fix a stochastic online scheduling instance
\[
I = (\mathcal{J}^*,\, \mathcal{M},\, \alpha \mid \beta \mid \gamma,\,
(\Omega, \mathcal{F}, \mathbb{P}),\, \mathcal{V})
\]
with latent operation set $\mathcal{O}^*$ and finite target job set
$\mathcal{J}^{\mathrm{target}} \subseteq \mathcal{J}^*$. The scheduling
POMDP induced by $I$ is the stochastic shortest-path POMDP
\[
\mathcal{P}_I = (\mathsf{S}_I,\; \mathsf{A}_I,\; \mathsf{O}_I,\; T_I,\; T_I^0,\;
Z_I,\; c_I)
\]
defined as follows.
\end{definition}

\paragraph{Preliminary: Processing requirements and progress rates.}

For each operation $o \in \mathcal{O}^*$ and machine $M_i \in \mathcal{M}$, let $p_{io}(\omega) \in (0, \infty]$ denote the realized processing time: the number of discrete time steps machine $M_i$ would require to complete $o$ from scratch in scenario $\omega$. The machine environment $\alpha$ determines these as follows:
\begin{itemize}
    \item $\alpha_1 = \epsilon$: $p_{1o}(\omega) = P_j(\omega)$, where $J_j$ is the parent
    job of $o$.
    \item $\alpha_1 = P$: $p_{io}(\omega) = P_j(\omega)$ for all $i$.
    \item $\alpha_1 = Q$: $p_{io}(\omega) = P_j(\omega) / S_i(\omega)$, where
    $S_i(\omega)$ is the speed of machine $M_i$.
    \item $\alpha_1 = R$: $p_{io}(\omega) = P_{io}(\omega)$, arbitrary.
    \item $\alpha_1 \in \{O, J, F\}$: $p_{io}(\omega) = P_{io}(\omega)$ on the designated
    machine for $o$; $p_{io}(\omega) = +\infty$ otherwise.
\end{itemize}

We adopt the transferable-progress convention: each operation $o$ has normalized total work content $1$, and machine $M_i$ contributes progress at rate
\[
\rho_{io}(\omega) := \frac{1}{p_{io}(\omega)} \in [0, 1).
\]
For $\alpha_1 \in \{P, Q\}$, this convention is canonical: the factorized form of processing times gives a natural machine-independent work content, and partial progress transfers across machines without ambiguity. For $\alpha_1 = R$, the bare unrelated-machines notation specifies only that $p_{io}$ is an arbitrary machine-operation-specific quantity; it does not, by itself, endow operations with a common notion of partial progress. The definition $\rho_{io} := 1/p_{io}$ is therefore an additional semantic convention that we impose: it interprets $p_{io}$ as the time $M_i$ would need to complete the entire operation, so that each time step on $M_i$ contributes a $1/p_{io}$ fraction of the total work, and this fractional progress is pooled across machines. This choice is an explicit formulation of the standard $R$-machine environment under preemption rules such that its conventions on partial progress is consistent with the $P$ and $Q$ environments and the transferable-progress convention. 


\paragraph{State space $\mathsf{S}_I$.}

A state is a tuple
\[
s = (\omega,\, t,\, B,\, Y)
\]
where:
\begin{itemize}
    \item $\omega \in \Omega$ is the latent scenario, encoding all realized primitive
    random variables (processing times, revelation times, release dates, due dates,
    etc.). It is drawn once at $t = 0$ and preserved identically by the transition
    kernel.
    \item $t \in \mathbb{N}_0$ is the current discrete time.
    \item $B : \mathcal{M} \to \mathcal{O}^* \cup \{\varnothing\}$ is the current
    machine-assignment map, where $B(M_i) = o$ means machine $M_i$ is currently
    processing operation $o$, and $B(M_i) = \varnothing$ means $M_i$ is idle.
    \item $Y : \mathcal{O}^* \times \mathcal{M} \to \mathbb{N}_0$ is the per-machine cumulative processing ledger, where $Y(o, M_i)$ records the total number of discrete time steps during which machine $M_i$ has processed operation $o$, up to and including time $t$.
\end{itemize}

The effective transferable progress of operation $o$ at state $s$ is
\[
Z(o;\, s) := \sum_{i=1}^{m} \rho_{io}(\omega)\, Y(o, M_i) = \sum_{i=1}^{m}
\frac{Y(o, M_i)}{p_{io}(\omega)}.
\]
Operation $o$ is complete at state $s$ if and only if $Z(o;\, s) \ge 1$. We write
$\mathcal{C}(s) := \{o \in \mathcal{O}^* : Z(o;\, s) \ge 1\}$ for the set of completed
operations at state $s$.

The state space is
\[
\mathsf{S}_I := \{(\omega, t, B, Y) \in \Omega \times \mathbb{N}_0 \times
(\mathcal{O}^* \cup \{\varnothing\})^{\mathcal{M}} \times
\mathbb{N}_0^{\mathcal{O}^* \times \mathcal{M}} :
(\omega, t, B, Y) \in \mathrm{Feas}_I\},
\]
where $(\omega, t, B, Y) \in \mathrm{Feas}_I$ if and only if all of the following hold:

\begin{itemize}
    \item[\textbf{(F1)}] \textbf{Revelation.} No unrevealed operation accrues
    processing:
    \[
    V_j(\omega) > t \implies Y(o, M_i) = 0 \quad \text{for all operations } o \text{ of }
    J_j \text{ and all } M_i.
    \]

    \item[\textbf{(F2)}] \textbf{Machine routing.} The assignment $B$ and ledger $Y$
    respect the machine constraints of $\alpha$:
    \begin{itemize}
        \item If $\alpha_1 = J$: $Y(O_{kj}, M_i) > 0$ only if $M_i = M_{\mu_{kj}}$, and
        $B(M_i) = O_{kj}$ only if $M_i = M_{\mu_{kj}}$.
        \item If $\alpha_1 = F$: $Y(O_{kj}, M_i) > 0$ only if $i = k$, and
        $B(M_i) = O_{kj}$ only if $i = k$.
        \item If $\alpha_1 = O$: $Y(O_{kj}, M_i) > 0$ only if $i = k$, and
        $B(M_i) = O_{kj}$ only if $i = k$.
    \end{itemize}

    \item[\textbf{(F3)}] \textbf{Single processing.} Each machine processes at most one
    operation: $|\{o \in \mathcal{O}^* : B(M_i) = o\}| \le 1$ for each $M_i$. Each
    operation is processed on at most one machine simultaneously:
    $|\{M_i : B(M_i) = o\}| \le 1$ for each $o$.

    \item[\textbf{(F4)}] \textbf{Non-preemption.} If $\beta_1 \neq \mathrm{pmtn}$: if
    $o \notin \mathcal{C}(s)$ and $Y(o, M_i) > 0$ for some $M_i$, then
    $B(M_j) \neq o$ and $Y(o, M_j) = 0$ for all $M_j \neq M_i$. That is, once an
    operation begins on a machine, it stays on that machine until completion.

    \item[\textbf{(F5)}] \textbf{Precedence.} If
    $\beta_3 \in \{\mathrm{prec}, \mathrm{tree}\}$: if $J_j \prec J_k$ in $G^*$ and
    any operation of $J_k$ has $B(M_i) = o_k$ or $Y(o_k, M_i) > 0$, then all
    operations of $J_j$ are complete.

    \item[\textbf{(F6)}] \textbf{Release dates.} If $\beta_4 = r_j$:
    $B(M_i) = o$ for an operation $o$ of $J_j$ only if $R_j(\omega) \le t$.

    \item[\textbf{(F7)}] \textbf{Resources.} If $\beta_2 = \mathrm{res}$: the
    simultaneously active operations $\{o : B(M_i) = o \text{ for some } i\}$ satisfy
    $\sum_{o\,\mathrm{active}} r_{ho}(\omega) \le \mathrm{cap}_h$ for each resource
    $R_h$.

    \item[\textbf{(F8)}] \textbf{Shop ordering.} If $\alpha_1 \in \{J, F\}$: if
    $B(M_i) = O_{k+1,j}$ or $Y(O_{k+1,j}, M_i) > 0$, then $O_{kj} \in \mathcal{C}(s)$.
\end{itemize}

We additionally define the following derived quantities at state
$s = (\omega, t, B, Y)$:
\begin{itemize}
    \item \textbf{Revealed job set:}
    $\mathcal{J}_t(\omega) := \{J_j \in \mathcal{J}^* : V_j(\omega) \le t\}$.
    \item \textbf{Revealed operation set:} $\mathcal{O}_t(\omega)$, the operations of
    jobs in $\mathcal{J}_t(\omega)$.
    \item \textbf{Revealed precedence sub-DAG}
    (if $\beta_3 \in \{\mathrm{prec}, \mathrm{tree}\}$):
    $G_t(\omega) := G^*[\mathcal{J}_t(\omega)]$. By assumption of consistency between revelation times and precedence constraints, this sub-DAG contains all predecessors of every
    revealed job.
    \item \textbf{Eligible operations:}
    $\mathcal{E}(s) \subseteq \mathcal{O}_t(\omega) \setminus \mathcal{C}(s)$,
    consisting of all revealed, incomplete operations satisfying all active constraints
    from $\alpha, \beta$ at state $s$.
\end{itemize}

\paragraph{Observation space $\mathsf{O}_I$ and observation kernel $Z_I$.}

The observation kernel is deterministic:
\[
Z_I(o \mid s) := \mathbf{1}[o = \mathrm{Obs}_I(s)],
\]
where $\mathrm{Obs}_I(s)$ for $s = (\omega, t, B, Y)$ produces the tuple
\[
o = \bigl(\mathcal{J}_t,\; G_t,\; B,\; Y|_{\mathcal{O}_t},\;
\mathcal{C}_t^{\mathrm{obs}},\; \Lambda^{\mathrm{prior}},\;
\Lambda_t^{\mathrm{rev}},\; \Lambda_t^{\mathrm{comp}}\bigr)
\]
with the following components.

\medskip
\noindent\textbf{$\Lambda^{\mathrm{prior}}$: Time-zero knowledge.} This component is
constant across all time steps and all $\omega$, encoding structural and distributional
information available to the scheduler from $t = 0$:
\begin{itemize}
    \item The machine environment $\alpha$, job environment $\beta$, and objective
    $\gamma$.
    \item The number of machines $m$ and their identities.
    \item Deterministic machine-level parameters (e.g., if $\alpha_1 = Q$ and machine
    speeds are deterministic, the values $S_i$).
    \item The prior distributional information for latent primitives: the joint law (or
    relevant marginals and moments) of processing times, revelation times, release
    dates, due dates, and any other random primitives under $\mathbb{P}$, to the extent
    specified by the observation convention of the instance.
\end{itemize}

\medskip
\noindent\textbf{$\Lambda_t^{\mathrm{rev}}$: Cumulative revelation-time information.}
For each job $J_j \in \mathcal{J}_t(\omega)$, the metadata revealed upon $J_j$'s arrival
at time $V_j(\omega)$:
\begin{itemize}
    \item The identity of $J_j$, its operation structure, and machine routing (if
    $\alpha_1 \in \{J, F, O\}$).
    \item Realized release date $R_j(\omega)$ (if $\beta_4 = r_j$).
    \item Realized due date $D_j(\omega)$ (if $\gamma$ depends on due dates and the
    convention specifies they are known upon revelation).
    \item Job weight $w_j$ (if relevant to $\gamma$).
    \item Distributional parameters or relevant moments of the unrealized processing
    requirements of $J_j$'s operations (e.g., $\mathbb{E}[P_j]$ in the
    Megow--Uetz--Vredeveld convention).
\end{itemize}
All precedence arcs among revealed jobs are already encoded in $G_t$.

\medskip
\noindent\textbf{$\mathcal{C}_t^{\mathrm{obs}}$: Observed completion set.} The set of
operations whose completion has been revealed to the scheduler by time $t$:
\[
\mathcal{C}_t^{\mathrm{obs}} := \{o \in \mathcal{O}_t(\omega) :
\text{completion of } o \text{ has been revealed by time } t\}.
\]
Concretely, $o \in \mathcal{C}_t^{\mathrm{obs}}$ if and only if $o \in \mathcal{C}(s)$
and the time step at which $Z(o;\,s)$ first reached $1$ has already passed. Since
completion is detected at the end of the time step in which cumulative progress crosses
the threshold, and the realized processing time is revealed at that moment (entering
$\Lambda_t^{\mathrm{comp}}$), we have $\mathcal{C}_t^{\mathrm{obs}} = \mathcal{C}(s)$
for all feasible states. Nevertheless, the distinction matters: $\mathcal{C}_t^{\mathrm{obs}}$
is an observed quantity derived from $\Lambda_t^{\mathrm{comp}}$, while $\mathcal{C}(s)$
is defined via the latent progress $Z(o;\,s)$ involving the unrevealed $p_{io}(\omega)$.
Admissibility and policies are stated in terms of $\mathcal{C}_t^{\mathrm{obs}}$, not
$\mathcal{C}(s)$.

\medskip
\noindent\textbf{$\Lambda_t^{\mathrm{comp}}$: Cumulative completion-time information.}
For each completed operation $o \in \mathcal{C}_t^{\mathrm{obs}}$, the realized
processing requirements on machines that contributed:
\[
\Lambda_t^{\mathrm{comp}} = \bigl\{(o,\, M_i,\, p_{io}(\omega)) :
o \in \mathcal{C}_t^{\mathrm{obs}},\; Y(o, M_i) > 0\bigr\}.
\]
This uses the cumulative convention: $\Lambda_t^{\mathrm{comp}}$ records all completions
up to time $t$. Since $\mathcal{C}_t^{\mathrm{obs}}$ is monotone in $t$, this is
equivalent to accumulating the per-step newly-completed operations.

\begin{remark}[Redundancy of $B$]
The assignment $B$ is included in the observation for convenience. Since the scheduler
knows its own past actions, it can reconstruct $B$ from the action history. Its
inclusion carries no extra exogenous information.
\end{remark}

\begin{remark}[Observation-visible vs.\ hidden information]
The observation $\mathrm{Obs}_I(s)$ reveals: which jobs exist and their metadata (for
revealed jobs only); the precedence structure among revealed jobs; the current assignment
and per-machine cumulative processing on revealed operations; the observed completion
set; realized processing times of completed operations; and distributional information
about unrealized primitives. It does \emph{not} reveal: the latent scenario $\omega$;
revelation times, identities, or existence of unrevealed jobs; realized processing times
of incomplete operations; or precedence arcs involving unrevealed jobs. In particular,
$\mathrm{Obs}_I$ is not injective: distinct latent scenarios producing identical visible
projections yield the same observation.
\end{remark}

The observation space $\mathsf{O}_I$ is the collection of all tuples of the above form
arising from feasible states.

\paragraph{Action space $\mathsf{A}_I$ and admissible actions.}

An action is a map
\[
a : \mathcal{M} \to \mathcal{O}^* \cup \{\varnothing\},
\]
with the interpretation that $a(M_i) = o$ assigns machine $M_i$ to process operation $o$
during the next time step, and $a(M_i) = \varnothing$ means $M_i$ idles.

The admissible action correspondence is defined on observations:
\[
\mathcal{U}_I : \mathsf{O}_I \rightrightarrows \mathsf{A}_I.
\]
An action $a \in \mathcal{U}_I(o)$ for observation
$o = (\mathcal{J}_t,\, G_t,\, B,\, Y|_{\mathcal{O}_t},\,
\mathcal{C}_t^{\mathrm{obs}},\, \Lambda^{\mathrm{prior}},\,
\Lambda_t^{\mathrm{rev}},\, \Lambda_t^{\mathrm{comp}})$ if and only if:
\begin{itemize}
    \item[\textbf{(A1)}] For each $M_i$ with $a(M_i) = o'$: the operation $o'$ is
    revealed ($o' \in \mathcal{O}_t$), incomplete
    ($o' \notin \mathcal{C}_t^{\mathrm{obs}}$), and eligible with respect to all
    visible constraints: precedence (via $G_t$), release dates (via
    $\Lambda_t^{\mathrm{rev}}$), shop ordering, and resources.

    \item[\textbf{(A2)}] Machine routing constraints from $\alpha$ are satisfied.

    \item[\textbf{(A3)}] Single processing: each operation is assigned to at most one
    machine.

    \item[\textbf{(A4)}] If $\beta_1 \neq \mathrm{pmtn}$: if $o'$ has accrued
    processing ($\sum_i Y(o', M_i) > 0$) and
    $o' \notin \mathcal{C}_t^{\mathrm{obs}}$, then $a(M_i) = o'$ only on the machine
    that has been processing $o'$.
\end{itemize}
Every quantity referenced in \textbf{(A1)--(A4)} is visible in the observation $o$. The
scheduler need not (and cannot!) reference the latent scenario $\omega$ to determine
admissibility.

\paragraph{Transition kernel $T_I$.}

For $s = (\omega, t, B, Y)$ and $a \in \mathcal{U}_I(\mathrm{Obs}_I(s))$, the transition
is the Dirac measure
\[
T_I(\cdot \mid s, a) := \delta_{\mathrm{Step}_I(s, a)},
\]
where
\[
\mathrm{Step}_I\bigl((\omega, t, B, Y),\; a\bigr) := (\omega,\; t+1,\; a,\; Y')
\]
with
\[
Y'(o, M_i) := Y(o, M_i) + \mathbf{1}[a(M_i) = o]
\]
for all $o \in \mathcal{O}^*,\; M_i \in \mathcal{M}$. Since admissibility enforces single
processing \textbf{(A3)}, this sum is either $0$ or $1$ for each $o$. The new assignment
is $B_{t+1} := a$. The latent scenario $\omega$ is preserved identically. Completions,
newly revealed jobs, and all other derived quantities update implicitly through the
derived-quantity definitions applied to the new state.

\begin{remark}[Conditional determinism]
All stochasticity in $\mathcal{P}_I$ is concentrated in the initial draw
$\omega \sim \mathbb{P}$. Conditional on $\omega$, the controlled process
$(t, B_t, Y_t)_{t \ge 0}$ is a deterministic dynamical system driven by the action
sequence. This is the hallmark of a latent-parameter POMDP: uncertainty is purely
epistemic!
\end{remark}

\paragraph{Initial state distribution $T_I^0$.}

The initial state is $s_0 = (\omega, 0, B_0, \mathbf{0})$ where
$\omega \sim \mathbb{P}$, $B_0 \equiv \varnothing$ (all machines idle), and $\mathbf{0}$
denotes the identically-zero ledger. Thus:
\[
T_I^0 := \mathbb{P} \otimes \delta_{(0,\, \varnothing,\, \mathbf{0})}.
\]

\paragraph{Cost function $c_I$ and objective.}

Define the terminal set
\[
\mathsf{S}_I^{\mathrm{term}} := \{s \in \mathsf{S}_I : \text{all operations of jobs in }
\mathcal{J}^{\mathrm{target}} \text{ lie in } \mathcal{C}(s)\},
\]
and the completed-job set
$\mathcal{C}_{\mathrm{jobs}}(s) := \{J_j : \text{all operations of } J_j \text{ lie in }
\mathcal{C}(s)\}$. The per-step cost depends on the scheduling objective $\gamma$:

\begin{itemize}
    \item \textbf{Makespan} ($\Gamma = \mathbb{E}[C_{\max}]$):
    \[
    c_I(s, a) := \begin{cases} 1 & \text{if } s \notin \mathsf{S}_I^{\mathrm{term}},
    \\ 0 & \text{if } s \in \mathsf{S}_I^{\mathrm{term}}. \end{cases}
    \]
    Then
    $\sum_{t=0}^{\infty} c_I(s_t, a_t) = \tau_{\mathrm{term}} = C_{\max}$, where
    $\tau_{\mathrm{term}} := \inf\{t \ge 0 : s_t \in \mathsf{S}_I^{\mathrm{term}}\}$.

    \item \textbf{Weighted completion time}
    ($\Gamma = \mathbb{E}[\sum_j w_j C_j]$):
    \[
    c_I(s, a) := \sum_{J_j \in \mathcal{J}^{\mathrm{target}} \setminus
    \mathcal{C}_{\mathrm{jobs}}(s)} w_j.
    \]
    Then $\sum_{t=0}^{\infty} c_I(s_t, a_t) = \sum_j w_j C_j$.

    \item Tardiness and unit-penalty objectives are encoded analogously with
    deadline-dependent per-step costs.
\end{itemize}

The expectation $\mathbb{E}^{\pi}$ integrates over the initial draw
$\omega \sim \mathbb{P}$ and the policy $\pi$. The expectation is not part of
the cost function; it arises from the stochastic shortest-path optimization criterion
applied to trajectories.

\paragraph{Restriction to special cases.}

\begin{remark}[Restriction to offline stochastic scheduling]
If $V_j(\omega) = 0$ for all $J_j \in \mathcal{J}^*$ and all $\omega \in \Omega$, then
$\mathcal{J}_t(\omega) = \mathcal{J}^*$ for all $t \ge 0$. The POMDP remains genuinely
partially observed: the full job and operation structure is known from $t = 0$ (via
$\Lambda^{\mathrm{prior}}$ and $\Lambda_0^{\mathrm{rev}}$), but realized processing
times remain hidden until completion. The revelation component
$\Lambda_t^{\mathrm{rev}}$ is constant for $t \ge 0$, and all new information arrives
through $\Lambda_t^{\mathrm{comp}}$.
\end{remark}

\begin{remark}[Restriction to online deterministic scheduling]
If all primitive random variables are almost surely constant (i.e.,
$|\mathrm{supp}(\mathbb{P})| = 1$ up to the revelation times), then the latent scenario
$\omega$ carries no hidden information beyond which jobs have not yet been revealed.
The observation $\mathrm{Obs}_I(s)$ determines the full state $s$ uniquely among
feasible states with the same revealed job set. The POMDP collapses to a fully-observed
MDP on the revealed state, with deterministic transitions. The remaining sequential
difficulty is that jobs arrive over time.
\end{remark}

\begin{remark}[Restriction to offline deterministic scheduling]
If additionally $V_j \equiv 0$, the full instance is known and deterministic from
$t = 0$. The POMDP collapses to a combinatorial optimization problem over feasible
schedules: no sequential decision-making under uncertainty remains.
\end{remark}

%%% ----------------------------------------------------------------
\subsubsection{Solutions and Policies}

We now define the scheduler's decision rules for the scheduling POMDP
$\mathcal{P}_I$. First, we more rigorously formalize the notion of nonanticipation via the observation history as touched on in the informal discussion of policies earlier.

\begin{definition}[Observation history]\label{def:obshist}
For each $t \ge 0$, the observation history up to time $t$ is a tuple
\[
h_t := (o_0,\, a_0,\, o_1,\, a_1,\, \ldots,\, a_{t-1},\, o_t)
\in \mathsf{H}_t^I,
\]
where $o_s \in \mathsf{O}_I$ and $a_s \in \mathsf{A}_I$ for all $s$. The space of
observation histories of length $t$ is
\[
\mathsf{H}_t^I := \begin{cases} \mathsf{O}_I & \text{if } t = 0, \\
(\mathsf{O}_I \times \mathsf{A}_I)^{t} \times \mathsf{O}_I & \text{if } t \ge 1.
\end{cases}
\]
We endow $\mathsf{H}_t^I$ with the product $\sigma$-algebra $\mathcal{H}_t^I$
generated by the coordinate maps.
\end{definition}

\begin{definition}[Randomized nonanticipative policy]\label{def:randpolicy}
A randomized nonanticipative policy for the scheduling POMDP $\mathcal{P}_I$
is a family of stochastic kernels
\[
\pi = (\pi_t)_{t \ge 0}, \qquad
\pi_t : \mathsf{H}_t^I \to \mathcal{P}(\mathsf{A}_I),
\]
such that for every $t \ge 0$ and every observation history $h_t \in \mathsf{H}_t^I$:
\begin{itemize}
    \item[\textbf{(P1)}] \textbf{Measurability.} The map
    $h_t \mapsto \pi_t(\cdot \mid h_t)$ is measurable with respect to
    $\mathcal{H}_t^I$.
    \item[\textbf{(P2)}] \textbf{Admissibility.} $\pi_t(\cdot \mid h_t)$ is supported
    on the admissible action set at the terminal observation:
    \[
    \pi_t\bigl(\mathcal{U}_I(o_t) \;\big|\; h_t\bigr) = 1,
    \]
    where $o_t$ is the final component of $h_t$.
\end{itemize}
The family of all such policies is denoted $\Pi_I$.
\end{definition}

\begin{remark}[Nonanticipation via the observation filtration]
The nonanticipation requirement is enforced structurally: at time $t$, the kernel
$\pi_t$ receives only the observation history
$h_t = (o_0, a_0, \ldots, a_{t-1}, o_t)$, which is the scheduler's complete
information record up to time $t$. Since the observation kernel $Z_I$ reveals only
$\mathrm{Obs}_I(s_t)$ (omitting unrevealed jobs, unrealized processing times, and
all other latent data) the policy cannot depend on any information not yet available
to the scheduler.

More precisely, the scheduling POMDP $\mathcal{P}_I$ together with a policy $\pi$
induces a probability measure $\mathbb{P}^{\pi}$ on the canonical trajectory space.
Under $\mathbb{P}^{\pi}$, define the observation filtration
\[
\mathcal{F}_t^{\mathrm{obs}} := \sigma(o_0, a_0, \ldots, a_{t-1}, o_t),
\qquad t \ge 0.
\]
Then the action $a_t$ selected by $\pi$ is $\mathcal{F}_t^{\mathrm{obs}}$-measurable
in the sense that $\pi_t(\cdot \mid h_t)$ depends on the trajectory only through
$\mathcal{F}_t^{\mathrm{obs}}$. This is the standard formalization of nonanticipation:
the controlled process $(a_t)_{t \ge 0}$ is adapted to
$(\mathcal{F}_t^{\mathrm{obs}})_{t \ge 0}$.
\end{remark}

\begin{remark}[Policies depend on observations, not latent states]
Since admissibility is defined on observations
($\mathcal{U}_I : \mathsf{O}_I \rightrightarrows \mathsf{A}_I$) and the policy kernel
conditions on observation histories, no component of the policy references the latent
scenario $\omega$. Two trajectories of latent states that produce the same observation
history receive the same action distribution.
\end{remark}

\begin{remark}[Belief-state sufficient statistic]\label{rmk:belief}
By standard POMDP theory, the posterior belief
\[
b_t := \mathrm{Law}(\omega \mid o_0, a_0, \ldots, a_{t-1}, o_t)
\in \mathcal{P}(\Omega)
\]
is a sufficient statistic of the observation history $h_t$ for optimal control. That is,
there exists an optimal policy that depends on $h_t$ only through $b_t$. Equivalently,
one may reformulate the scheduling POMDP as a fully-observed MDP on the belief space
$\mathcal{P}(\Omega)$, the belief MDP, in which the state at time $t$ is $b_t$
and the transition is given by Bayesian updating. The optimal cost-to-go is then a
function of $b_t$ alone. This connects directly to the posterior
$\Pi_t = \mathrm{Law}(\theta \mid \mathcal{H}_t)$ appearing in the centralized ITAP state defined in the next chapter, and is the structural reason a belief-augmented state space is the natural formulation for that model.
\end{remark}

\begin{definition}[Deterministic nonanticipative policy]\label{def:detpolicy}
A deterministic nonanticipative policy is a family of measurable maps
\[
\varphi = (\varphi_t)_{t \ge 0}, \qquad \varphi_t : \mathsf{H}_t^I \to \mathsf{A}_I,
\]
such that $\varphi_t(h_t) \in \mathcal{U}_I(o_t)$ for every $h_t \in \mathsf{H}_t^I$.
\end{definition}

\begin{remark}
A deterministic nonanticipative policy is a special case of a randomized nonanticipative
policy, obtained by taking $\pi_t(\cdot \mid h_t) := \delta_{\varphi_t(h_t)}$ for each
$t$ and $h_t$.
\end{remark}

\begin{definition}[Policy value and the scheduling optimization problem]
\label{def:policyvalue}
For a policy $\pi \in \Pi_I$, define the expected cost
\[
J_I(\pi) := \mathbb{E}^{\pi}\!\left[\sum_{t=0}^{\infty} c_I(s_t, a_t)\right],
\]
where $\mathbb{E}^{\pi}$ denotes expectation under the measure $\mathbb{P}^{\pi}$
induced by the initial draw $\omega \sim \mathbb{P}$, the transition kernel $T_I$, the
observation kernel $Z_I$, and the action selection under $\pi$.

The scheduling optimization problem is
\[
J_I^* := \inf_{\pi \in \Pi_I} J_I(\pi).
\]
For the makespan cost ($c_I \equiv 1$ until termination), this reduces to
$J_I^* = \inf_\pi \mathbb{E}^{\pi}[C_{\max}]$, recovering the stochastic scheduling
objective $\Gamma = \mathbb{E}[C_{\max}]$. Analogous reductions hold for weighted
completion time, tardiness, and unit-penalty objectives via the corresponding per-step
costs defined previously.
\end{definition}

\subsection{Examples}

\subsubsection{Offline deterministic: \(P\mid \mathrm{prec}\mid C_{\max}\)}

Consider the identical-parallel-machine makespan problem with precedence constraints. Here, we have $n$ jobs with processing times $p_j$ and precedence constraints given by a DAG $G$, and the goal is to schedule the jobs on $m$ identical machines to minimize the makespan $C_{\max}$.

Graham’s list scheduling algorithm gives a \((2-\tfrac1m)\)-approximation for this problem on \(m\) identical machines. The problem is NP-hard in general; even special cases such as \(P2\mid\mathrm{chains}\mid C_{\max}\) are weakly NP-hard, and \(P2\mid \mathrm{prec},,p_j\in{1,2}\mid C_{\max}\) is $W[2]$-hard when parameterized by partial-order width. On the positive side, for unit-size jobs on a fixed number of machines, a quasi-polynomial time approximation scheme (QPTAS) is known, and for \(m=3\), the exact complexity of the unit-size case remains famously delicate.

A nice way to present the algorithm is this. Fix any linear extension \(L\) of the precedence order. Whenever a machine becomes idle, schedule the first available job in \(L\). Two canonical lower bounds on the optimal makespan are

\[
\frac{1}{m}\sum_{j} p_j
\qquad\text{and}\qquad
\ell_{\max},
\]
where \(\ell_{\max}\) is the length of a longest precedence chain. If \(j\) is the last job completed by list scheduling, then the interval before \(j\) starts can be decomposed into busy periods and forced-idle periods. Busy periods contribute at most \(\frac{1}{m}\sum_i p_i-\frac{p_j}{m}\), while forced-idle periods can be charged to a chain ending at \(j\), contributing at most \(\ell_{\max}-p_j\). Hence
\[
C_{\max}^{LS}\le \frac1m\sum_i p_i+\Bigl(1-\frac1m\Bigr)\ell_{\max}
\le \Bigl(2-\frac1m\Bigr)\mathrm{OPT}.
\]

\subsubsection{Offline stochastic: \(P\mid \mathrm{prec}\mid \mathbb E[C_{\max}]\), and the weighted-completion extension}

For the stochastic offline case, consider the stochastic counterpart of the same model as before. Here, Graham’s list scheduling bound extends directly from deterministic makespan to expected makespan, namely, $P\mid \mathrm{prec}\mid \mathbb E[C_{\max}]$ still admits a \((2-\tfrac1m)\)-approximation. We do note, however, Skutella and Uetz show that once one moves from makespan to expected weighted completion time with precedence constraints, the problem becomes much subtler: ordinary Graham-style list scheduling can be arbitrarily bad, and LP-based delayed list scheduling is needed to recover constant-factor guarantees. Their paper gives the first constant-factor approximation algorithms for stochastic parallel-machine scheduling with precedence constraints.

The clean proof sketch for the makespan version is pathwise. For each realization \(\omega\), let
\[
W(\omega)=\sum_j P_j(\omega),\qquad
H(\omega)=\text{length of a longest realized precedence chain}.
\]
Running the same list-scheduling rule as above gives, for each \(\omega\),
\[
C_{\max}^{LS}(\omega)\le \frac{W(\omega)}{m}+\Bigl(1-\frac1m\Bigr)H(\omega).
\]
Since every feasible schedule must satisfy
\[
\mathrm{OPT}(\omega)\ge \max\Bigl\{\frac{W(\omega)}{m},,H(\omega)\Bigr\},
\]
taking expectations yields
\[
\mathbb E[C_{\max}^{LS}]
\le \Bigl(2-\frac1m\Bigr)\mathbb E[\mathrm{OPT}],
\]
and the stochastic problem is NP-hard already because the deterministic case is its degenerate special case.

In a related note, for
\[
P\mid \mathrm{prec}\mid \mathbb E\left[\sum_j w_j C_j\right],
\]
Skutella and Uetz combine an LP relaxation with delayed list scheduling to obtain constant-factor approximations; reporting a guarantee of
\[
3+2\sqrt2-\frac{1+\sqrt2}{m}
\]
for the precedence-constrained case. Indeed, once the objective changes from makespan to weighted completion, stochastic scheduling rapidly stops being a straightforward extension of the deterministic theory.

\subsubsection{Stochastic online: classical SOS for \(P\mid r_j\mid \mathbb E[\sum_j w_j C_j]\)}

Consider a problem where jobs arrive over time, on arrival the scheduler learns the job’s weight and expected processing time, but the realized processing time remains unknown until completion. In this model, the number of jobs is not known in advance, jobs arrive online in release-date order, and the objective is expected weighted completion time. When release dates are absent on a single machine, weighted shortest expected processing time (WSEPT) is the natural and well-known rule, namely, order jobs based on the highest ratio of weight to expected processing time. However, with release dates, we can consider the \(\alpha\)-Shift-WSEPT policy. Moreover, on parallel machines, one may combine it with fixed-assignment policies such as Modified MinIncrease or RandAssign.

On one machine, \(\alpha\)-Shift-WSEPT modifies each release date to
\[
r'_j=\max{r_j,\alpha,\mathbb E[P_j]}
\]
and then always runs the available job of maximum Smith ratio \(w_j/\mathbb E[P_j]\). Megow, Uetz, and Vredeveld prove that this is a \((\delta+2)\)-approximation for
\[
1\mid r_j\mid \mathbb E\left[\sum_j w_j C_j\right]
\]
when processing times are \(\delta\)-NBUE, with the best choice \(\alpha=1\). For the parallel-machine version
\[
P\mid r_j\mid \mathbb E\left[\sum_j w_j C_j\right]
\]

This modified MinIncrease policy assigns each arriving job to the machine minimizing an incremental expected-cost expression and then sequences jobs on that machine by \(\alpha\)-Shift-WSEPT. This gives the bound
\[
\rho = 1+\max \left\{1+\frac{\delta}{\alpha}, \alpha+\delta+\frac{(m-1)(\Delta+1)}{2m}\right\},
\]
where \(\Delta\) bounds the squared coefficient of variation. For NBUE processing times, optimizing \(\alpha\) gives a ratio below \(3.62-\frac{1}{2m}\). Their randomized policy RandAssign, which assigns each revealed job uniformly at random to a machine and then runs \(\alpha\)-Shift-WSEPT locally, achieves the same bound.

This example also gives a clean online lower-bound. Because deterministic processing times are contained as a special case, classical deterministic online lower bounds transfer immediately to SOS. Megow, Uetz, and Vredeveld recall that in the deterministic single-machine online-time model, a \(2\)-competitive algorithm exists and this lower bound is tight; for identical parallel machines, they cite a lower bound of \(1.309\) for deterministic online algorithms, a best-known deterministic upper bound of \(2.62\), and a randomized upper bound of \(2\). They also prove lower bounds internal to their SOS framework: any fixed-assignment policy has asymptotic lower bound about \(1.24\), and their MinIncrease policy has lower bound \(3/2\) for general distributions. 


\subsection{Theorem Proving as a Scheduling Problem}

Consider